% Generated draft: requires mathematical and source review.
\documentclass[11pt]{article}
\usepackage[margin=25mm]{geometry}
\usepackage{fontspec}
\setmainfont{Latin Modern Roman}
\usepackage{amsmath,amssymb,mathtools,unicode-math}
\setmathfont{Latin Modern Math}
\usepackage{enumitem,needspace,xurl,hyperref}
\setlength{\parindent}{0pt}
\setlength{\parskip}{5pt}
\setlength{\emergencystretch}{3em}
\newcommand{\N}{\mathbb{N}}
\newcommand{\Z}{\mathbb{Z}}
\newcommand{\Q}{\mathbb{Q}}
\newcommand{\R}{\mathbb{R}}
\newcommand{\C}{\mathbb{C}}
\newcommand{\F}{\mathbb{F}}
\newcommand{\A}{\mathbb{A}}
\newcommand{\PP}{\mathbb{P}}
\newcommand{\E}{\mathbb{E}}
\newcommand{\eps}{\varepsilon}
\newcommand{\dd}{\,\mathrm{d}}
\begin{document}
\section*{UnsolvedMath ID: 30001286}
\textbf{Source problem code:} catalogue entry
\subsection{Definitions and mathematical statement}
Write \(\mathcal H=HI(k)\), \(\mathcal H_n=HI^{\leq n}(k)\), and let \(I_n:\mathcal H_n\hookrightarrow\mathcal H\) denote the inclusion appearing in the supplied record. The record associates \(\mathcal H_n\) with \(n\)-motivic sheaves, but does not define either category, specify the hypotheses on \(k\), give the range of \(n\), or state the topology, coefficient ring, or transfer conventions. These are unresolved prerequisites to a fully specified statement; the notation here retains the source's unspecified choices.

The principal question is whether \(I_n\) is left exact. Under the categorical convention that left exactness means preservation of finite limits, this asks whether, for every finite diagram \(D:J\to\mathcal H_n\), the comparison morphism
\[
I_n\!\left(\varprojlim_J D\right)\longrightarrow\varprojlim_J(I_n\circ D)
\]
is an isomorphism, assuming the relevant finite limits exist. The existence of these limits and the intended exactness convention must be checked against the source definitions.

The original statement additionally asks whether the inclusion admits a left adjoint
\[
L_n=(-)^{\leq n}:\mathcal H\longrightarrow\mathcal H_n.
\]
Precisely, this asks for a functor \(L_n\) and bijections
\[
\operatorname{Hom}_{\mathcal H_n}(L_nF,G)\cong\operatorname{Hom}_{\mathcal H}(F,I_nG)
\]
natural in every \(F\in\mathcal H\) and \(G\in\mathcal H_n\). The supplied original statement presents this as a stronger question. That source-specific characterization is retained, but its relationship to the first question requires the missing categorical hypotheses and conventions.
\subsection{Short English statement}
The problem asks whether the stated inclusion of \(n\)-motivic sheaves preserves finite limits, and additionally whether it has a left adjoint. The definitions of the categories, the hypotheses on \(k\), and the range of \(n\) are missing from the supplied record, so this is a conditional formulation requiring source clarification and human mathematical review.
\subsection{Sources}
\begin{itemize}
\item Local problem registry: \texttt{data/unsolvedmath/problems.json}, permanent UnsolvedMath ID \(30001286\), source problem code \texttt{catalogue entry}, title \emph{Left Exactness of n-Motivic Sheaf Inclusions}. The supplied record gives the proposed year as 2009 and names no individual proposer.
\item Supplied report citation: \emph{Algebraic \(K\)-Theory and Motivic Cohomology} (2009), Oberwolfach Reports, pp.\ 1731--1774, DOI 10.4171/owr/2009/31. \url{https://doi.org/10.4171/owr/2009/31}. The background attributes the problem to workshop open-problem discussions documented in this report.
\item Supplied literature reference: Ayoub--Barbieri-Viale on \(1\)-motivic sheaves. \url{https://arxiv.org/abs/math/0607738}.
\item Supplied literature assessment dated \texttt{2026-08-21}: the record labels the problem open and reports that the Oberwolfach report explicitly leaves left exactness open and states a stronger conjecture for \(n\)-motivic sheaves. This is supplied dated curation, not an independently verified literature assessment; the background requests independent verification of time-sensitive status and cited claims.
\item Supplied verification metadata: \texttt{statement\_status} is \texttt{corrected\_verified}; the accompanying text reports comparison with the official Oberwolfach report and normalization of notation, context, and question form without changing the source-level claim. This is a supplied verification assertion, not verification performed for this draft.
\end{itemize}
\end{document}
